\chapter{Introduction}
\underline{\textbf{\textit{SyncFlow}}} is a full-stack, web-based \textbf{Project Schedule Management and Operational Control platform} built around the six core processes of the Project Management Institute's \textbf{PMBOK\textsuperscript{\textregistered} Guide}: \textit{Plan Schedule Management}, \textit{Define Activities}, \textit{Sequence Activities}, \textit{Estimate Activity Durations}, \textit{Develop Schedule}, and \textit{Control Schedule}. Designed as a \textit{shared, role-based web workspace}, it brings formal scheduling mathematics, day-to-day task tracking, and Natural Language Processing together inside a \textbf{zero-cost open-source stack}.

\vspace{4mm}

\textit{SyncFlow} was created to fix a practical gap in project management software, one that becomes especially visible in technical fields like software engineering, video game development, building construction, or design of laboratory equipment, in short any project where activities are highly interdependent. In practice, teams usually have to settle for tools at one of two extremes. \textbf{Lightweight task boards} like \textit{Trello} and \textit{Asana} make daily coordination easy, yet offer little in the way of \textbf{Operations Research}: they cannot catch dependency loops, compute \textit{Critical Path Method (CPM)} dates and float margins, or estimate \textit{probabilistic schedule risk}. Conversely, \textbf{enterprise desktop suites} like \textit{Microsoft Project} and \textit{Oracle Primavera P6} support formal scheduling calculations, but come with steep learning curves, high licensing fees, \textit{desktop-bound workflows} kept away from frontline contributors, and no built-in language tools to help draft an initial plan.

\vspace{4mm}

SyncFlow tackles both sides of this problem with a \textbf{dual-engine web architecture} pairing exact scheduling algorithms with \textbf{human-in-the-loop Natural Language Processing}. At the core, a deterministic Python engine runs \textbf{topological sorting}, \textbf{cycle detection}, \textbf{CPM} forward and backward passes, \textbf{Three-Point PERT} probability estimates, \textbf{Earned Value Management} (\textbf{EVM}) S-Curves, and automated \textbf{14-point schedule audits} based on Defense Contract Management Agency (\textbf{DCMA}) criteria. Alongside this math engine, an NLP module powered by \textbf{Google Gemini Flash} speeds up initial drafting by turning plain-text scope documents into \textbf{Work Breakdown Structures} (\textbf{WBS}), proposing PERT duration ranges, and extracting \textit{precedence relationships} from workflow descriptions. The application is primarily designed for Project Managers, while other users (Project Members) have more freedom, but are provided with a more restricted set of functionalities. 

\vspace{4mm}

Just as importantly, the platform keeps schedule plans connected to actual work through \textbf{separate, synchronized interfaces} for different roles. \textbf{Project Managers} set up \textit{governance policies}, \textit{working calendars}, \textit{dependency networks}, and \textit{baselines}, while \textbf{Team Members} work from a focused \textit{task agenda} with \textbf{live stopwatch timers} and \textbf{deliverable evidence submissions}. Whenever a contributor logs work hours or finishes a verified activity, project-wide Earned Value metrics and critical path buffers update in \textbf{real time}, keeping the schedule grounded in real progress rather than manual estimates.
